% Header reference level
\clearpairofpagestyles
\chead{\headmark}
\automark[section]{section}
\ifoot{\ifootertext}
\ofoot{\ofootertext}
\cfoot{\thepage}

\section{Conclusion}
\label{sec:Conclusion}

This codebook provides operational guidance for the planning and execution of encoding institutions using the Institutional Grammar 2.0. This information is supplementary to the conceptual introduction provided in \citet{Frantz2021,Frantz2022InstitutionalGrammar}. 

Following a general introduction of the principles and motivations of the Institutional Grammar, the codebook initially outlines the theoretical concepts underlying IG 2.0, including the coding on multiple levels of expressiveness (IG Core, IG Extended, IG Logico) in \cref{sec:Definitions}, before supporting the navigation through the codebook based on the reader's objectives and background in the form of a Reader's Guide in \cref{subsec:ReadersGuide}. This is followed by the discussion of essential document preparation steps (pre-coding steps) in \cref{sec:PrecodingSteps} that complement the broader design considerations outlined in \citet{Frantz2022InstitutionalGrammar}, central aspects of which are captured in the Study Design Checklist appended to this codebook (see Appendix \ref{sec:Checklist}).

Following the general methodological orientation, the codebook provides detailed coding guidelines for regulative and constitutive statements across all IG 2.0 levels of expressiveness (\cref{sec:CodingGuidelines}). This is followed by advanced concepts, such as the discussion of encoding hybrid institutional statements (\cref{subsec:ConstitutiveRegulativeHybrids}), i.e., institutional statements consisting of both regulative and constitutive components, polymorphic institutional statements (\cref{subsec:SyntacticPolymorphs}), and the discussion of structural patterns (\cref{subsec:DataStructurePatterns}). 

A mechanism to document the parameterization of IG 2.0 studies based on the selective application of IG 2.0 components and features -- in the form of IG 2.0 Profiles -- is provided in \cref{subsec:IgProfiles}. The listing of taxonomies to support the semantic annotation of institutional information in the previous section (\cref{sec:Taxonomies}) concludes the substantive part of the codebook.

It is important to note that the guidelines provided in this codebook are of general nature and emphasize the \begin{emp}operational use\end{emp} of IG 2.0, with specific focus on the encoding of institutional information. Doing so, they may not capture specifics potentially relevant for a given project, let alone the detailed introduction of the underlying theoretical concepts. For an in-depth introduction into the Institutional Grammar, existing applications, challenges that motivate the IG 2.0 and detailed conceptual foundations, the reader is referred to companion literature~\citep{Frantz2021,Frantz2022InstitutionalGrammar}. 

However, the codebook aims to support the development of supplementary project-specific guidelines that consider application context (e.g., domain, language, types of documents, legal traditions, etc.) and analytical objectives (i.e., evaluation of encoded statements) more explicitly, aspects that a general codebook cannot consider. As indicated above, the reader's guide in \cref{subsec:ReadersGuide} intends to highlight relevant sections. It is furthermore important to note that the coding instructions provided here are intentionally tool-agnostic, and open to adaptation to arbitrary encoding means (e.g., spreadsheets, text annotation tools, etc.), which themselves may be augmented with specific guidance. 

For a complete overview of relevant resources, including a theoretical treatment of the underlying concepts and principles, as well as resources that support operational coding (e.g., cheat sheet, tool-specific guidance, software), please refer to \url{https://newinstitutionalgrammar.org/resources}. Given the ongoing development and application efforts, the website is updated whenever relevant new material becomes available.

Please further note that these guidelines will be continuously refined based on theoretical developments, feedback from users as well as ongoing empirical validation efforts. To retrace subsequent changes, please note the specific version and version history of these guidelines outlined at the beginning of this document. Irrespective of refinements, all revisions of the codebook will be retained for future reference.

%\newpage
